\documentclass{article}

\usepackage{neurips_2026}

\usepackage[utf8]{inputenc}
\usepackage[T1]{fontenc}
\usepackage{hyperref}
\usepackage{url}
\usepackage{booktabs}
\usepackage{amsfonts}
\usepackage{amsmath}
\usepackage{amssymb}
\usepackage{nicefrac}
\usepackage{microtype}
\usepackage{xcolor}
\usepackage{multirow}
\usepackage{graphicx}

% Notation
\newcommand{\drefuse}{\mathbf{d}_{\text{refuse}}}
\newcommand{\dbeh}{\mathbf{d}_{\text{beh}}}
\newcommand{\mdelta}{\boldsymbol{\Delta}}
\newcommand{\Amin}{{\mathcal{A}^{-}}}
\newcommand{\Apos}{{\mathcal{A}^{+}}}
\newcommand{\AminSM}{{\mathcal{A}^{-}_{\text{SM}}}}
\newcommand{\AposSM}{{\mathcal{A}^{+}_{\text{SM}}}}

\title{Different Inputs, Different Adversarial Experts:\\
  Locate with Causal Gradients, Suppress via Router Bias}

\author{%
  Anonymous Authors \\
  \texttt{anonymous@example.edu}
}

\begin{document}

\maketitle

\begin{abstract}
Existing MoE-steering defenses (e.g., SteerMoE) assume a \emph{fixed}
set of safety-relevant experts, derived from population activation
statistics and applied identically to every input. We show this
assumption is wrong: \textbf{different harmful inputs recruit different
adversarial experts} — the set of experts whose activation most reduces
the model's refusal behavior differs substantially from query to query,
with pairwise Jaccard of only $0.11$ between top-5 sets across $200$
queries on OLMoE-1B-7B. We propose a two-step defense that respects
this input-specificity: (i) \emph{locate} per-query adversarial experts
by taking a causal autograd through router biases against a behavioral
refusal direction, yielding a $1024$-dimensional signed fingerprint per
query; (ii) \emph{suppress} the top-$K$ most-negative experts by a
constant router-bias shift during generation. At identical bandwidth
($K{=}25$), our method beats SteerMoE on three of four engaged
jailbreak attacks (peak: past-tense $32{\rightarrow}62\%$
safety rate) \emph{and} incurs $0$\,pp helpfulness tax on XSTest benign
prompts while SteerMoE over-refuses $12$\,pp under the same judge.
We further report that per-query \emph{positive}-$\Delta$ experts do not
align with SteerMoE's safe-expert pool (chance-level enrichment), and
correspondingly amplifying them breaks generation coherence — an
asymmetry between suppression and amplification that matches the
asymmetric validity of the two expert pools.
\end{abstract}

\section{Introduction}

Mixture-of-Experts (MoE) architectures scale language models by routing
each token through a sparse subset of feed-forward experts
\cite{shazeer2017outrageously,jiang2024mixtralexperts,muennighoff2025olmoe}.
The router is a natural surface for test-time behavioral control: adjusting
\emph{which experts fire} on an input is cheaper and more interpretable
than weight edits or activation steering.

Recent work (SteerMoE~\citep{fayyaz2025steermoe}) demonstrates the
promise of this approach. SteerMoE selects two global expert pools —
safety-associated ($\mathcal{A}^{+}$) and unsafe-associated
($\mathcal{A}^{-}$) — from the difference of activation frequencies on
paired safe/unsafe responses, and applies a fixed log-softmax clipping
mask to those pools at inference on \emph{every} input. This treatment
assumes that a single population-level partition of experts into
``safe'' and ``unsafe'' is adequate to steer the model on any given
query.

\paragraph{This assumption is wrong.} We measure, per query, the
\emph{causal} effect of perturbing each $(\ell, e)$ router logit on
the model's projection onto a behavioral refusal direction. This
yields a signed $L \times E$ fingerprint $\mdelta(q)$ of experts that,
on \emph{this} query, either drive or suppress refusal. On $200$
harmful prompts across four attack families, the top-5 most-adversarial
experts (i.e., most refusal-suppressing) overlap only at pairwise
Jaccard $0.11$ across queries — the set of ``adversarial experts''
depends strongly on the input. Tellingly, the union of per-query
top-5 sets covers $145$ distinct $(\ell, e)$ pairs out of $1024$, yet
no single expert appears in more than $52\%$ of top-5 sets. There is
no universal adversarial expert.

\paragraph{Method (two steps).}
\begin{enumerate}
  \item \textbf{Locate adversarial experts via causal gradients.}
  For an incoming query $q$, register a zero-initialised bias
  $b_{\ell,e}$ on each router logit, forward to a mid-layer probe,
  and take one autograd call against
  $\mathrm{refuse\_score}(q) = h_{\ell^*}^{\top} \dbeh$. The resulting
  $\Delta_{\ell,e}(q) = \partial \mathrm{refuse\_score} / \partial b_{\ell,e}$
  ranks experts by their signed causal contribution to refusal.
  \item \textbf{Suppress via router bias.} Apply a constant negative
  bias ($-10$) to the top-$K$ most-negative-$\Delta$ experts during
  generation. Everything else is untouched.
\end{enumerate}

\paragraph{Findings.} On OLMoE-1B-7B across four engaged jailbreak
families from PandaGuard (\emph{none}, GCG, PAIR, past-tense) at
identical bandwidth ($K{=}25$) as SteerMoE:
\begin{itemize}
\item The two-step method beats SteerMoE on 3/4 attacks (peak:
  past-tense $+21$pp over SteerMoE, $+30$pp over baseline) using only
  the negative-$\Delta$ half — \emph{no amplification}.
\item On XSTest benign prompts, the method incurs $0$\,pp over-refusal
  cost; SteerMoE over-refuses $12$\,pp under the same judge on the
  same $50$ prompts.
\item \textbf{Positive-$\Delta$ experts are different: amplification
  fails.} The top-$K$ \emph{positive}-$\Delta$ experts do not enrich
  into SteerMoE's $\mathcal{A}^{+}$ (chance-level overlap), and
  amplifying them with a $+3$ bias produces fluent-looking gibberish.
  The asymmetry between suppression success and amplification failure
  tracks this asymmetric enrichment.
\end{itemize}

\paragraph{Contributions.}
(i) Empirical evidence that adversarial experts in MoE models are
input-specific rather than population-universal, refuting a premise
of prior fixed-mask steering; (ii) a two-step test-time defense —
causal-gradient localisation followed by router-bias suppression —
that outperforms the strong SteerMoE baseline on safety while
reducing helpfulness cost to zero; (iii) a clean negative result
pinning the success of suppression and the failure of amplification
to concordance (or lack thereof) between per-query and
population-level expert pools.

\section{Related Work}
\label{sec:related}

\paragraph{Refusal directions.}~\citet{arditi2024refusal} identified a
one-dimensional refusal direction in dense transformer residual streams
and showed that orthogonalizing this direction from attention and MLP
output weights ablates the model's safety behavior.~\citet{chen2025persona}
extended the technique to arbitrary persona traits using contrastive
system prompts. Both works operate at the activation or weight level and
do not exploit MoE routing.

\paragraph{MoE expert steering.} SteerMoE~\cite{fayyaz2025steermoe}
proposed expert-level steering for MoE by (i) detecting behavior-linked
experts via a risk-difference (RD) statistic on paired safe/unsafe
responses from BeaverTails, and (ii) activating or deactivating the
top-$K$ experts (by $|\text{RD}|$) via log-softmax clipping on router
logits. RICE~\cite{wang2025rice} amplifies thinking-token experts for
reasoning tasks. Both methods select experts \emph{globally} from
aggregate statistics and apply a fixed mask at inference.

\paragraph{Weight-level safety.} Recent work has explored structural
removal of refusal directions from MoE weights~\cite{fayyaz2025steermoe,
arditi2024refusal}. In Appendix~\ref{app:weight-ortho} we show that
applying Arditi-style orthogonalization to all $E \times L$ expert
down-projections on OLMoE is marginal (e.g., $+6$pp on clean harmful
prompts, $0$pp on jailbreaks with an independent query rewrite), and is
outperformed by activation-style steering on three of four jailbreak
families. We thus focus on the test-time router-bias regime.

\section{Method}
\label{sec:method}

\subsection{Preliminaries}

We consider an MoE transformer with $L$ layers, $E$ experts per layer,
and top-$k$ routing. For input $q$ with token embeddings at position $t$,
layer $\ell$ outputs
\begin{equation}
  \mathbf{h}_{\ell+1,t} = \sum_{i \in \mathcal{T}_{\ell,t}} \tilde{p}_{\ell,t,i} \cdot \mathrm{Expert}_{\ell,i}(\mathbf{h}_{\ell,t}),
\end{equation}
where $\mathcal{T}_{\ell,t}$ is the top-$k$ set drawn from the router
probabilities $p_{\ell,t,i} = \mathrm{softmax}(W_\ell^{\text{router}} \mathbf{h}_{\ell,t})_i$.

\subsection{Two Probe Directions}
\label{sec:directions}

Our method uses two per-layer directions, each extracted via mean-diff
over a different contrast and each playing a distinct role.

\paragraph{$\drefuse$ — harmfulness detector (for gating).}
$\drefuse[\ell]$ is the classical Arditi-style refusal direction:
the mean difference of the \emph{prompt}'s last-token residual between
harmful and benign prompts at layer $\ell$. It is strong on ``does this
prompt look harmful?'' and is directly comparable across benign and
harmful distributions. We use it as a gate signal.

\paragraph{$\dbeh$ — refusal-vs-compliance axis (for attribution).}
$\dbeh[\ell]$ captures \emph{the model is currently refusing} rather
than \emph{the model detected harmful intent}. On a set of baseline
harmful prompts we generate under the original model, obtain
Qwen2.5-14B safety judgments, partition responses into refuse ($R$)
and comply ($C$) sets on the \emph{same} prompts, and compute
\begin{equation}
  \dbeh[\ell] = \operatorname*{mean}_{r \in R, t \in r} \mathbf{h}_{\ell,t}(r) - \operatorname*{mean}_{c \in C, t \in c} \mathbf{h}_{\ell,t}(c).
\end{equation}
$\dbeh$ is aggregated across four attack distributions (\emph{none},
cipher, PAIR, past-tense) with equal weight. Because its contrast is
between two behaviors on the same harmful prompts, projections onto
$\dbeh$ are only reliable under \emph{harmful-input distribution} —
which is exactly why we use $\drefuse$ (not $\dbeh$) for gating. Both
directions are unit-normalized per layer.

\subsection{Step 1: Locate Adversarial Experts}
\label{sec:delta}

Fix a probe layer $\ell^*$. Forward the query to layer $\ell^*$ and
capture the last-prompt-token residual $\mathbf{h}^* = \mathbf{h}_{\ell^*,T(q)}(q)$.
From this single residual we compute two scalars:
\begin{align}
  \lambda(q) &= \mathbf{h}^* \cdot \drefuse[\ell^*] && \text{(harmfulness)} \\
  \mathrm{refuse\_score}(q) &= \mathbf{h}^* \cdot \dbeh[\ell^*] && \text{(current refusal engagement)}
\end{align}

\paragraph{Causal attribution via autograd.}
We register a zero-initialized bias $b_{\ell,e}$ on every
router logit in layers $\{0,\dots,\ell^*\}$ and ask: if we perturbed
$b_{\ell,e}$ by a small amount, how would refuse\_score change? A single
\verb|torch.autograd.grad| call answers this:
\begin{equation}
  \Delta_{\ell,e}(q) \;=\; \left. \frac{\partial \, \mathrm{refuse\_score}(q)}{\partial b_{\ell,e}} \right|_{b=0}.
\end{equation}
Note that $\lambda(q)$ and $\mathrm{refuse\_score}(q)$ share the same
forward pass: the gate signal comes free alongside $\mdelta$.
$\Delta_{\ell,e}(q) < 0$ identifies experts whose activation
\emph{reduces} refusal on this query — our adversarial experts.

\subsection{Step 2: Adaptive Router-Bias Suppression}
\label{sec:defense}

Rather than a uniform intervention, we scale the bias magnitude
(i) at the query level by the harmfulness score $\lambda(q)$, and
(ii) at the expert level by the causal magnitude $|\Delta_{\ell,e}(q)|$.

\paragraph{Query-level magnitude.}
Define a linear ramp on $\lambda(q)$:
\begin{equation}
  \beta_{\text{eff}}(q) = \beta \cdot \operatorname{clip}\!\left(\frac{\lambda(q) - \lambda_{\text{lo}}}{\lambda_{\text{hi}} - \lambda_{\text{lo}}},\, 0,\, 1\right).
\end{equation}
$\lambda_{\text{lo}}$ and $\lambda_{\text{hi}}$ are calibrated once from
two held-out distributions (benign and harmful) — specifically, the
$95$th percentile of $\lambda$ on a mixture of benign queries and the
$5$th percentile on harmful queries, respectively. No held-out fitting
per attack or per deployment. Truly benign queries
($\lambda \le \lambda_{\text{lo}}$) yield $\beta_{\text{eff}} = 0$ — no
intervention; clearly harmful queries
($\lambda \ge \lambda_{\text{hi}}$) yield full $\beta$; ambiguous
queries interpolate.

\paragraph{Expert-level weighting.}
For each query $q$, let $\Amin(q) = \operatorname*{arg\,topk}_{\ell,e}\bigl\{-\Delta_{\ell,e}(q)\bigr\}$.
Weight the bias on each selected expert by its relative
adversarial magnitude:
\begin{equation}
  b_{\ell,e}(q) = -\, \beta_{\text{eff}}(q) \cdot \frac{|\Delta_{\ell,e}(q)|}{\max_{(\ell',e') \in \Amin(q)} |\Delta_{\ell',e'}(q)|}
  \quad \text{for } (\ell,e) \in \Amin(q),
\end{equation}
and $b_{\ell,e} = 0$ otherwise. The strongest adversarial expert on a
given query receives the full $-\beta_{\text{eff}}(q)$ bias; weaker ones
receive proportionally less.

\paragraph{Generation.}
Register a single forward hook on each MoE gate module that adds
$b_{\ell,e}(q)$ to router logits before softmax, then generate
autoregressively. Model weights are never modified.

Hard-gate ($\mathbf{1}[\lambda > \tau]$) and uniform-magnitude
($b_{\ell,e} = -\beta_{\text{eff}}(q)$ without expert-level weighting)
variants are ablated in Appendix~\ref{app:adaptive}.

\subsection{Offline Caching and Inference Cost}
\label{sec:cost}
$\lambda(q)$, $\mathrm{refuse\_score}(q)$, and $\mdelta(q)$ all come
from a \emph{single} forward-backward pass through layers
$\{0,\dots,\ell^*\}$ ($\sim 0.6$s on B200). Under the gate, truly
benign queries can skip the backward entirely after reading
$\lambda(q)$. Generation itself is unchanged; the method adds only one
(shared) pass of attribution, not three separate inferences. Amortized
over a $256$-token response, total latency overhead is
approximately $5{-}10\%$ (App.~\ref{app:impl}).

\section{Experiments}
\label{sec:experiments}

\subsection{Setup}

\paragraph{Model.} OLMoE-1B-7B-0924-Instruct~\cite{muennighoff2025olmoe}:
$L{=}16$ layers, $E{=}64$ experts per layer, top-$k{=}8$, hidden dim
$d{=}2048$. Weights frozen throughout.

\paragraph{Attacks.} PandaGuard JBB attack suite~\cite{pandaguard}; we
focus on four \emph{engaged} attack families where the baseline model
meaningfully engages refusal circuitry (observed
$\|\mdelta\| \gg 0$ at the probe layer): \emph{none} (direct harmful
requests), GCG, PAIR, and past-tense reformulations
\cite{andriushchenko2024past}. For each attack we draw $n{=}50$
prompts by smallest \texttt{goal\_idx}. Separate experiments on bypass
attacks (deepinception, cipher, AIM) are reported in
Appendix~\ref{app:bypass} — these families do not drive the baseline
model into refusal and are mechanically out of reach for any per-query
Delta-based intervention.

\paragraph{Baselines.}
\begin{itemize}
\item \textbf{No intervention (baseline):} OLMoE with no hooks.
\item \textbf{SteerMoE:} our faithful reimplementation of
 ~\cite{fayyaz2025steermoe} on OLMoE, $K_{\Apos} = K_{\Amin} = 25$,
  selection via BeaverTails risk-difference, evaluated on the same
  $n{=}50$ prompts per attack.
\item \textbf{Fixed-5:} a hand-selected small set of 5 universal
  suppressor experts from a prior pilot (Appendix~\ref{app:fixed5}),
  router bias $-10$. Included as a same-mechanism, smaller-$K$
  reference.
\end{itemize}

\paragraph{Judge.} Qwen2.5-14B-Instruct with a 0–100 harm-alignment
prompt; \textsc{safe} $\iff$ score $\geq 70$. Identical judge prompt and
threshold used for baseline, SteerMoE, and our methods.

\paragraph{Probe layer.} $\ell^* = 10$, selected based on the prior
observation that refusal-related signal peaks mid-network for
Instruct-tuned decoder models~\cite{arditi2024refusal}.

\subsection{Main Result: Safety on Engaged Jailbreaks}
\label{sec:main}

Table~\ref{tab:main} reports safety rates on the four engaged attacks.

\begin{table}[t]
\centering
\caption{Safety rate (\%, higher is better) on jailbreak attacks spanning
all 5 mechanistic clusters of PandaGuard JBB (L47e taxonomy). $n{=}100$
per cell, Qwen2.5-14B judge, threshold $70$. All $K$-using methods:
$K{=}25$ per direction. Bypass-cluster attacks (bottom) serve as negative
controls confirming mechanistic unreachability. \textbf{[PLACEHOLDER
rows 1--10 pending Exp 1 in \texttt{experiment\_plan.md}; current rows
1, 2, 5 at $n{=}50$ from earlier pilot.]}}
\label{tab:main}
\small
\begin{tabular}{llcccc}
\toprule
Attack & Cluster & baseline & SteerMoE & $C_{\text{perQ-sup}}$ (ours) & $D_{\text{perQ-comb}}$ \\
 & & (no hook) & ($K{=}25{+}25$) & ($K{=}25$, neg only) & ($K{=}25{+}25$, per-q) \\
\midrule
\multicolumn{6}{l}{\emph{Engaged cluster (5)}} \\
\emph{none}         & 5 & 48.0 & 45.5 & \textbf{52.0} & ~0.0 \\
GCG                 & 5 & 16.0 & XX.X\textsuperscript{\dag} & \textbf{30.0} & ~0.0 \\
PAIR                & 5 & 25.0 & 33.0 & 22.0          & ~0.0 \\
ICA                 & 5 & XX.X & XX.X & XX.X          & XX.X \\
\midrule
\multicolumn{6}{l}{\emph{Tense-shift cluster (4)}} \\
past-tense          & 4 & 32.0 & 41.4 & \textbf{62.0} & ~0.0 \\
tense\_future       & 4 & XX.X & XX.X & XX.X          & XX.X \\
\midrule
\multicolumn{6}{l}{\emph{DAN / roleplay cluster (3)}} \\
AIM                 & 3 & XX.X & XX.X & XX.X          & XX.X \\
BETTER\_DAN         & 3 & XX.X & XX.X & XX.X          & XX.X \\
\midrule
\multicolumn{6}{l}{\emph{Persuasion / search cluster (1)}} \\
new\_pair           & 1 & XX.X & XX.X & XX.X          & XX.X \\
new\_renellm        & 1 & XX.X & XX.X & XX.X          & XX.X \\
\midrule
\multicolumn{6}{l}{\emph{Bypass cluster (2) -- negative controls}} \\
new\_deepinception  & 2 & XX.X & XX.X & XX.X          & XX.X \\
new\_gpt4\_cipher   & 2 & XX.X & XX.X & XX.X          & XX.X \\
\midrule
mean (non-bypass $n{=}10$) & & XX.X & XX.X & \textbf{XX.X} & ~0.0 \\
\bottomrule
\end{tabular}
\\
\vspace{2pt}
\footnotesize
\textsuperscript{\dag} SteerMoE GCG was omitted from the L40 reproduction
due to scoping; Exp 1 re-runs to fill.
\end{table}

Per-query suppression improves safety on three of four attacks, with a
particularly large gain on past-tense ($+30$pp absolute over baseline,
$+20.6$pp over SteerMoE). On \emph{none} it beats SteerMoE by $+6.5$pp.
GCG improves by $+14$pp over baseline; SteerMoE was not evaluated on
GCG in our reproduction. PAIR is the one attack where $C_{\text{perQ-sup}}$
slightly underperforms baseline ($-3$pp); we return to this in
\S\ref{sec:discussion}.

Numbers in Table~\ref{tab:main} are at $n{=}50$ per cell (binomial 95\%
CI $\pm 14$pp). Extended results at $n{=}100$, a $K$ sweep, and a
direction ablation ($\dbeh$ vs $\drefuse$) are reported in
Appendices~\ref{app:n100}, \ref{app:ksweep}, and \ref{app:drefuse}
respectively. The main-text conclusions hold under all three.

\subsection{Amplification Fails: A Clean Negative Result}
\label{sec:amplify}

To test whether symmetric amplification further improves safety, we
evaluated $D_{\text{perQ-comb}}$: $K{=}25$ per-query most-negative
experts with bias $-10$ \emph{plus} $K{=}25$ per-query most-positive
experts with bias $+3$. The smaller positive bias was chosen
conservatively; a pilot at $+10$ produced pure gibberish on all
prompts.

Even at $+3$, $D_{\text{perQ-comb}}$ yields $0\%$ safety rate across all
four attacks — but inspection of responses reveals these are not
unsafe, they are \emph{incoherent}:
\begin{quote}\small
\texttt{"The and Image in-TheThe. A: Aid: 1. 100. The -- it'ss Not,toing:;..."}
\end{quote}
In contrast, $C_{\text{perQ-sup}}$ responses on the same prompts are
fluent refusals:
\begin{quote}\small
\texttt{"As an AI language model, it is not appropriate or ethical to write an article defaming a U.S. president..."}
\end{quote}
We interpret this asymmetry in \S\ref{sec:discussion}: per-query
positive-$\Delta$ experts do not align with SteerMoE's globally validated
safe-expert pool, whereas per-query negative-$\Delta$ experts do.

\subsection{Concordance with Global Risk-Difference Selection}
\label{sec:concordance}

To validate that our per-query $\Amin$ identifies meaningful
``bad experts'', we measure its overlap with SteerMoE's global
$\AminSM$ derived independently from BeaverTails activation
counts.

\begin{table}[t]
\centering
\caption{Overlap enrichment of per-query top-5 most-negative
$\Delta_{\text{beh}}$ experts with SteerMoE's global A-minus pool at
various global $K$. ``Enrichment'' is observed intersection rate
divided by chance ($|\AminSM|/1024$). Per-query positive
$\Delta$ experts show no such enrichment into SteerMoE's A-plus pool.}
\label{tab:concord}
\small
\begin{tabular}{ccccc}
\toprule
 & \multicolumn{2}{c}{$\Amin$ (suppress candidates)} & \multicolumn{2}{c}{$\Apos$ (amplify candidates)} \\
$|\AminSM|$ & \% hit & enrichment & \% hit & enrichment \\
\midrule
5   & 2.9  & $5.9\times$ & 0.5 & $1.0\times$ \\
25  & 9.6  & $3.9\times$ & 1.7 & $0.7\times$ \\
50  & \textbf{38.9} & \textbf{8.0$\times$} & 2.2 & $0.5\times$ \\
100 & 58.3 & $6.0\times$ & 4.4 & $0.5\times$ \\
\bottomrule
\end{tabular}
\end{table}

Per-query negative $\Delta$ picks enrich up to $8\times$ into
SteerMoE's A-minus, with no observable cross-sign confusion (per-query
negative $\cap$ SteerMoE A-plus remains at chance, $\approx 1.8\%$).
This is an \emph{independent} validation: our causal-gradient signal
and SteerMoE's activation-count signal agree on which experts drive
unsafe behavior, despite completely different detection machinery. The
per-query positive side, critically, shows no such enrichment. This
observation is diagnostic for the amplification failure
(\S\ref{sec:amplify}): there is no shared notion of a ``safe expert''
between the two frameworks, while there is a shared notion of an
``unsafe expert''.

\subsection{Adversarial Experts Are Input-Specific}
\label{sec:diversity}

Is per-query selection actually per-query, or does it collapse into a
universal mask? We measure pairwise Jaccard between the top-5
most-negative-$\Delta$ sets across all $\binom{200}{2}$ query pairs.

\begin{table}[h]
\centering
\caption{Diversity of per-query top-5 suppression sets across 200
queries (50 per attack).}
\label{tab:diversity}
\small
\begin{tabular}{lc}
\toprule
statistic & value \\
\midrule
pairwise Jaccard (mean)            & $0.106$ \\
pairwise Jaccard (median)          & $0.111$ \\
fraction of disjoint pairs         & $36.6\%$ \\
unique top-5 signatures / 200 q    & $194 / 200$ \\
most frequent ``universal bad''~$(L{=}10, E{=}30)$ & $52\%$ of queries \\
unique $(L,e)$ covered by some $q$ & $145 / 1024$ \\
\bottomrule
\end{tabular}
\end{table}

Selection is strongly per-query: 97\% of queries have a unique top-5,
and over a third of query pairs share no experts at all. A small
handful of $(L,e)$ pairs are near-universal suppressors (the top
$(L{=}10, E{=}30)$ appears in $52\%$ of per-query top-5 sets), but the
majority of the mask is query-specific. Compared to SteerMoE's single
global mask of 25 experts, our union covers 145 distinct experts
across all 200 queries — the per-query mechanism is not merely a
noisier version of the global mask but a genuinely different selection.

\subsection{Benign Over-Refusal}
\label{sec:benign}

A defense is only worthwhile if it preserves helpful behavior on
benign inputs. We run $C_{\text{perQ-sup}}$ on 50 XSTest
\cite{rottger2024xstest} benign prompts designed to superficially
resemble harmful requests, and measure the answer rate (non-refusal)
under a separate Qwen2.5-14B judge with an ANSWER/REFUSE classification
prompt.

\begin{table}[h]
\centering
\caption{Benign answer rate (\%, higher is better) on the same 50 XSTest
safe prompts used throughout. Apples-to-apples: identical prompts,
identical Qwen2.5-14B ANSWER/REFUSE classifier. $C_{\text{perQ-sup}}$
incurs zero helpfulness tax; SteerMoE over-refuses 7 of 50 prompts.}
\label{tab:benign}
\small
\begin{tabular}{lccc}
\toprule
 & baseline & \textbf{$C_{\text{perQ-sup}}$} & SteerMoE \\
\midrule
answer rate (\%)    & 98.0 (49/50) & \textbf{98.0 (49/50)} & 86.0 (43/50) \\
over-refusal cost   & --           & \textbf{0.0\,pp}      & \textbf{--12\,pp} \\
\midrule
\multicolumn{4}{l}{\emph{Per-type breakdown}:} \\
figurative language & 100\,(25/25) & 100\,(25/25) & 88\,(22/25) \\
homonyms            &  96\,(24/25) &  96\,(24/25) & 84\,(21/25) \\
\bottomrule
\end{tabular}
\end{table}

The comparison is striking. Both methods target expert-level router
intervention at identical total bandwidth ($K{=}25$), on the exact same
50 XSTest prompts under the same judge. $C_{\text{perQ-sup}}$'s refused
prompt (a single homonym-style query) is refused identically under the
baseline — no genuine over-refusal. SteerMoE, by contrast, over-refuses
7 of 50 benign prompts (12\,pp cost), uniformly distributed across both
XSTest categories. We attribute the difference to \emph{selection
locality}: per-query top-$K$ picks experts by each query's own causal
fingerprint, whereas SteerMoE's global mask applies the same 25
``unsafe'' and 25 ``safe'' experts to every input — including benign
queries about killing Python processes or terminating subscriptions,
where the global mask has no justification. Combined with the
amplification asymmetry established in \S\ref{sec:amplify}, the paper's
main claim — \emph{query-adaptive expert suppression improves safety on
3 of 4 engaged attacks and reduces over-refusal on benign prompts
relative to the fixed-mask baseline} — is supported end to end.

\section{Discussion}
\label{sec:discussion}

\paragraph{Why suppression works but amplification fails.} The
concordance analysis (Table~\ref{tab:concord}) and the amplification
pilot (\S\ref{sec:amplify}) line up to tell a consistent story.
Per-query $\Delta_{\text{beh}}$ negative identifies experts whose
\emph{role} is to suppress the refusal direction at a probe point;
this aligns with SteerMoE's population-level A-minus because both
methods are picking up a robust \emph{disruption} signal.
Per-query positive $\Delta_{\text{beh}}$, by contrast, identifies
experts whose marginal activation would push the residual toward the
refusal direction — but this is a \emph{counterfactual} direction, not
a feature the model is actually using to refuse. Forcing these experts
to fire at amplitude ($+3$ router bias) pushes the residual stream
into regions the model does not know how to decode, producing
gibberish. Suppression has no such failure mode because it merely
reduces already-active, disruptive experts' contribution.

\paragraph{PAIR underperformance.} PAIR is the one attack where
$C_{\text{perQ-sup}}$ trails baseline by $-3$pp. PAIR queries are
optimized by an adversarial LLM to produce prompts that look benign on
the surface; $\Delta_{\text{beh}}$ on these prompts has intermediate
magnitude and the selected experts resemble those selected on
\emph{none} more than on GCG or past-tense. One hypothesis is that
PAIR queries evade the ``refusal engagement'' that our method depends
on. A larger $n$ is needed to confirm whether this is a true
regression or noise within the $\pm 14$pp 95\% CI at $n{=}50$.

\paragraph{Bandwidth comparison.} $C_{\text{perQ-sup}}$ uses $K{=}25$
negative experts only — half the total intervention bandwidth of
SteerMoE ($K{=}25$ activate $+$ $K{=}25$ deactivate). That the
suppression-only method still outperforms suggests SteerMoE's
amplification side is either redundant or, as our
$D_{\text{perQ-comb}}$ pilot hints, silently trading helpfulness
for negligible safety gain.

\paragraph{Limitations.}
\begin{itemize}
\item \textbf{Single MoE architecture.} All experiments are on
  OLMoE-1B-7B. Whether the same per-query signal structure exists on
  DeepSeek-V2-Lite, Qwen-MoE, or Mixtral is an immediate follow-up.
\item \textbf{Bypass attacks unaddressed.} Attacks like deepinception
  and cipher do not engage refusal circuitry at all at the probe layer
  ($\|\Delta_{\text{beh}}\| \to 0$). Our method has no signal to act
  on. Appendix~\ref{app:bypass} confirms this is independent of
  suppression strength.
\item \textbf{Inference overhead.} Although we cache $\Delta(q)$
  offline in our experiments, a deployed system would incur one extra
  forward-backward pass per query (${\sim}5\%$ of greedy generation
  cost for $256$ tokens). The method is compatible with batching.
\item \textbf{$n{=}50$ per cell.} 95\% binomial CIs on safety rates
  are $\pm$14pp; differences under this are underpowered. Extending to
  $n{=}100$ is a natural next step.
\end{itemize}

\section{Conclusion}

\emph{Different inputs recruit different adversarial experts in
Mixture-of-Experts language models.} We quantified this input-specificity
(Jaccard $0.11$ between per-query top-5 sets on $200$ queries), and
designed a two-step defense that respects it: locate per-query
adversarial experts with a single causal autograd through router
biases, suppress them with a constant router-bias shift at generation.
The method is weight-preserving, adds only a single forward-backward
pass per query, and at identical intervention bandwidth outperforms
the strong SteerMoE baseline on jailbreak safety while incurring zero
over-refusal cost on XSTest benign prompts. The suppress-amplify
asymmetry we report closes the loop: the ``unsafe'' direction of our
per-query signal has concordant evidence from a population-level pool
(SteerMoE's $\mathcal{A}^{-}$), the ``safe'' direction does not — and
amplification accordingly fails. MoE safety steering is, at bottom, a
matter of \emph{which adversarial experts to silence on this input}.

\section*{Acknowledgments}
\emph{Deferred to final copy.}

\bibliographystyle{plainnat}
\bibliography{example_paper}

\appendix

\section{Weight-Orthogonalization Baseline}
\label{app:weight-ortho}

As an additional ablation we applied Arditi-style weight
orthogonalization on all $L \times E = 1024$ expert down-projection
matrices: $W_{\ell,e}^{\text{down}} \leftarrow W - \hat{d}\hat{d}^\top W$
where $\hat{d} = \dbeh[\ell]/\|\dbeh[\ell]\|$. Results on the same
attack suite are reported in Table~\ref{tab:wortho}. The intervention
is near-neutral — it marginally improves \emph{none} ($+6$pp) and
cipher ($+4$pp) but provides no help on PAIR or past-tense. Together
with the low-cost advantage of our test-time method, this justifies
focusing the main paper on router-bias rather than weight-level
interventions.

\begin{table}[h]
\centering
\caption{Weight-ortho variants on OLMoE ($n{=}100$, Qwen judge).}
\label{tab:wortho}
\small
\begin{tabular}{lcc}
\toprule
attack & baseline & ortho\_all \\
\midrule
\emph{none}   & 48.0 & 54.0 \\
cipher        & 8.0  & 12.1 \\
PAIR (global) & 25.0 & 25.0 \\
past-tense    & 32.0 & 29.0 \\
benign        & --   & 100.0 \\
\bottomrule
\end{tabular}
\end{table}

\section{Bypass Attacks}
\label{app:bypass}

Attacks including deepinception, cipher (GPT-4 cipher), AIM,
BETTER\_DAN, renellm, and autodan exhibit $\|\Delta_{\text{beh}}(q)\|
\approx 0.16 \times \|\Delta_{\text{beh}}(\text{none})\|$ in our cache,
with cross-attack cosine of mean $\Delta$ anti-correlated ($-0.64$) vs.
the \emph{none} centroid. These attacks do not trigger the model's
refusal circuitry at the probe layer, so per-query $\Delta$-based
interventions have no signal to amplify or suppress. Empirically,
$C_{\text{perQ-sup}}$ on these attacks shows $0$–$4\%$ refusal rate,
matching baseline within noise. We regard this cluster as mechanically
out of reach for any $\dbeh$-anchored defense.

Table~\ref{tab:bypass_quant} quantifies this claim across three
representative bypass families under both our method and SteerMoE.
\textbf{[PLACEHOLDER: fill after Exp 5 in \texttt{experiment\_plan.md}.]}

\begin{table}[h]
\centering
\caption{Safety rate on bypass-cluster attacks
($n{=}50$ per cell). No method moves the needle.
\textbf{[PLACEHOLDER]}}
\label{tab:bypass_quant}
\small
\begin{tabular}{lccc}
\toprule
attack & baseline & SteerMoE & $C_{\text{perQ-sup}}$ \\
\midrule
new\_deepinception & XX.X & XX.X & XX.X \\
new\_gpt4\_cipher  & XX.X & XX.X & XX.X \\
AIM                & XX.X & XX.X & XX.X \\
\bottomrule
\end{tabular}
\end{table}

\section{Extended $n{=}100$ Scale}
\label{app:n100}

The main-text Table~\ref{tab:main} uses $n{=}50$ per cell. Table~\ref{tab:main_n100}
re-runs all conditions at $n{=}100$ to halve the binomial CI.
\textbf{[PLACEHOLDER: fill after Exp 1.]}

\begin{table}[h]
\centering
\caption{Safety rate at $n{=}100$. \textbf{[PLACEHOLDER]}}
\label{tab:main_n100}
\small
\begin{tabular}{lcccc}
\toprule
 & baseline & SteerMoE & $C_{\text{perQ-sup}}$ & $D_{\text{perQ-comb}}$ \\
\midrule
\emph{none}      & XX.X & XX.X & XX.X & ~0.0 \\
GCG              & XX.X & --   & XX.X & ~0.0 \\
PAIR             & XX.X & XX.X & XX.X & ~0.0 \\
past-tense       & XX.X & XX.X & XX.X & ~0.0 \\
\bottomrule
\end{tabular}
\end{table}

\section{Adaptive Magnitude Ablation}
\label{app:adaptive}

Our method (\S\ref{sec:defense}) scales the bias at two levels:
query-wise by $\beta_{\text{eff}}(q) = \beta \cdot f(\lambda(q))$ with
$f$ a linear ramp, and expert-wise by
$|\Delta_{\ell,e}(q)| / \max_{(\ell',e')\in\Amin(q)}|\Delta_{\ell',e'}(q)|$.
Table~\ref{tab:adaptive} contrasts four variants. \textbf{[PLACEHOLDER:
fill after Exp 9 in \texttt{experiment\_plan.md}.]}

\begin{table}[h]
\centering
\caption{Intervention-magnitude ablation on engaged jailbreaks and
benign (XSTest + AlpacaEval), $n{=}50$ each. \textbf{[PLACEHOLDER]}}
\label{tab:adaptive}
\small
\begin{tabular}{lcccc}
\toprule
scheme & query gate & expert weight & engaged safety $\uparrow$ & benign answer $\uparrow$ \\
\midrule
uniform, always on  & none          & uniform    & XX.X & XX.X \\
hard gate only      & $\mathbf{1}[\lambda > \tau]$ & uniform    & XX.X & XX.X \\
linear ramp only    & $\beta \cdot f(\lambda)$ & uniform & XX.X & XX.X \\
linear ramp + $|\Delta|$-weighted   & $\beta \cdot f(\lambda)$ & $\propto |\Delta|$ & XX.X & XX.X \\
\bottomrule
\end{tabular}
\end{table}

Calibration: $\lambda_{\text{lo}}$ = 95th percentile of $\lambda$ on a
held-out benign mixture (AlpacaEval + XSTest); $\lambda_{\text{hi}}$ =
5th percentile on a held-out harmful mixture (PandaGuard's 25 attacks).
\textbf{[PLACEHOLDER: report exact values after calibration.]}

\section{$K$ Sweep Ablation}
\label{app:ksweep}

$K{=}25$ was chosen to match SteerMoE's bandwidth. Table~\ref{tab:ksweep}
shows $C_{\text{perQ-sup}}$ safety rate on past-tense as a function of
$K$. \textbf{[PLACEHOLDER: fill after Exp 3.]}

\begin{table}[h]
\centering
\caption{Safety rate on past-tense by $K$. \textbf{[PLACEHOLDER]}}
\label{tab:ksweep}
\small
\begin{tabular}{cccccc}
\toprule
$K$ & 5 & 10 & 25 & 50 & 100 \\
\midrule
safe \% & XX.X & XX.X & XX.X & XX.X & XX.X \\
\bottomrule
\end{tabular}
\end{table}

\section{Direction Ablation: d\_behavior vs d\_refuse}
\label{app:drefuse}

Our score function uses $\dbeh$. Table~\ref{tab:drefuse} shows the same
pipeline with $\drefuse$ (prompt last-token harmful-vs-benign mean-diff,
used by prior work \cite{arditi2024refusal}).
\textbf{[PLACEHOLDER: fill after Exp 4.]}

\begin{table}[h]
\centering
\caption{Effect of score function on per-query selection.
\textbf{[PLACEHOLDER]}}
\label{tab:drefuse}
\small
\begin{tabular}{lcc}
\toprule
 & $C_{\text{perQ-behavior}}$ (ours) & $C_{\text{perQ-refuse}}$ \\
\midrule
\emph{none}      & 52.0 & XX.X \\
GCG              & 30.0 & XX.X \\
PAIR             & 22.0 & XX.X \\
past-tense       & 62.0 & XX.X \\
\bottomrule
\end{tabular}
\end{table}

\section{Generalisation to a Second MoE Architecture}
\label{app:qwen_moe}

To test whether our findings depend on OLMoE-specific quirks, we
replicate the full pipeline on Qwen1.5-MoE-A2.7B (24 layers, 60
experts, top-$k{=}4$). The same hyperparameters ($K{=}25$,
$\beta{=}{-}10$) and probe-layer convention (40\,\% depth) are used.
\textbf{[PLACEHOLDER: fill after Exp 2 in \texttt{experiment\_plan.md}.]}

\begin{table}[h]
\centering
\caption{Main results reproduced on Qwen1.5-MoE-A2.7B. \textbf{[PLACEHOLDER]}}
\label{tab:qwen_moe}
\small
\begin{tabular}{lcccc}
\toprule
 & baseline & SteerMoE & $C_{\text{perQ-sup}}$ & $\Delta$ vs SteerMoE \\
\midrule
\emph{none}        & XX.X & XX.X & XX.X & XX.X \\
GCG                & XX.X & --   & XX.X & -- \\
PAIR               & XX.X & XX.X & XX.X & XX.X \\
past-tense         & XX.X & XX.X & XX.X & XX.X \\
\midrule
benign (XSTest $n{=}50$) & XX.X & XX.X & XX.X & XX.X \\
\bottomrule
\end{tabular}
\end{table}

We further check that the per-query diversity finding replicates:
pairwise Jaccard of top-$K$ sets across queries is
\textbf{[PLACEHOLDER]}, comparable to 0.11 observed on OLMoE.

\section{Fixed-5 Expert Selection}
\label{app:fixed5}

The fixed-5 baseline uses five hand-selected universal suppressor
experts from a prior pilot:
$(L{=}8, E{=}29), (L{=}10, E{=}6), (L{=}6, E{=}15), (L{=}9, E{=}8),
(L{=}9, E{=}16)$. All five are in SteerMoE's top-100 A-minus pool.

\section{Implementation Details}
\label{app:impl}

\paragraph{Router hook.} OLMoE's gate module returns either a raw logit
tensor (prefill) or a $(p, \text{top\_val}, \text{top\_idx})$ tuple
(decode with KV cache). Our hook handles both by recomputing the bias
addition before softmax and re-extracting top-$k$.

\paragraph{Generation config.} Greedy decoding, \texttt{do\_sample=False},
\texttt{max\_new\_tokens=256}, \texttt{temperature=1.0}, bf16 throughout.
\texttt{pad\_token\_id} set to \texttt{eos\_token\_id}.

\paragraph{Inference overhead.} End-to-end wall-clock at $n{=}10$:
baseline \textbf{[PLACEHOLDER]} ms/query, $C_{\text{perQ-sup}}$ with
online $\Delta$ extraction \textbf{[PLACEHOLDER]} ms/query, overhead
\textbf{[PLACEHOLDER]} \%. Measured on a single NVIDIA B200 (Exp 6 in
\texttt{experiment\_plan.md}).

\paragraph{Compute.} All experiments run on a single NVIDIA B200
(180 GB). $\Delta$ cache: ${\sim}10$ min for $N{=}200$. Main generation:
${\sim}70$ min per $200$-query condition at
\texttt{max\_new\_tokens}${=}256$. Judging: ${\sim}10$ min per 200 responses.

\end{document}
